\chapterbackmatter{Solutions to Supplementary Exercises}

\begin{enumerate}
  \SupplementarySolution{1}{Counting a General Real Superfield}
  Reality makes \(C,M,N,D,V_\mu\) real: there are four scalar components
  plus four vector components, hence eight bosonic components.  Each of
  \(\omega,\lambda\) is Majorana and has four real components, hence eight
  fermionic components.  No equation has been used.  In particular,
  \(V_\mu\) is not yet transverse and \(M,N,D,\lambda\) are not eliminated;
  the equality is therefore an off-shell statement.

  \SupplementarySolution{2}{Closure on the Wess--Zumino Scalar}
  From \(\delta A=\bar\alpha\psi\),
  \[
  \begin{aligned}
    [\delta_\alpha,\delta_\beta]A
     ={}&\bar\beta\!
       \left[\partial_\mu(A+i\gamma_5B)\gamma^\mu\alpha
             +(F-i\gamma_5G)\alpha\right]-(\alpha\leftrightarrow\beta).
  \end{aligned}
  \]
  Majorana interchange identities make the \(F,G\) terms and the
  \(\partial B\) terms cancel, while the vector bilinear is antisymmetric.
  Thus
  \[
    [\delta_\alpha,\delta_\beta]A
      =2(\bar\beta\gamma^\mu\alpha)\partial_\mu A.
  \]
  Direct substitution gives the same translation on
  \(B,\psi,F,G\).  The auxiliaries are precisely what permits this closure
  without imposing a field equation on \(\psi\).

  \SupplementarySolution{3}{The Product Rule for Lowest Components}
  Multiplication of the two finite expansions, followed by the Majorana
  interchange rules, gives
  \[
  \begin{aligned}
    C&=C_1C_2,&
    \omega&=C_1\omega_2+C_2\omega_1,\\
    M&=C_1M_2+C_2M_1+\frac i2\bar\omega_1\gamma_5\omega_2,\\
    N&=C_1N_2+C_2N_1-\frac12\bar\omega_1\omega_2,\\
    V^\mu&=C_1V_2^\mu+C_2V_1^\mu
       -\frac i2\bar\omega_1\gamma_5\gamma^\mu\omega_2 .
  \end{aligned}
  \]
  These are Eqs.~\eqref{eq:26.2.19}--\eqref{eq:26.2.23}.  The signs in
  the last three lines arise when a fermionic component is moved through
  \(\bar\theta\).

  \SupplementarySolution{4}{Why the Covariant Superderivative Works}
  Write \(\SuperQ=A+B\) and \(\SuperD=A-B\), where \(A\) differentiates
  with respect to \(\theta\) and \(B\) is linear in \(\theta\partial_\mu\).
  Then \(\{A,A\}=\{B,B\}=0\), while the two cross terms cancel in
  \(\{\SuperD,\SuperQ\}\) and add in
  \(\{\SuperD,\bar{\SuperD}\}\), giving
  \[
    \{\SuperD_\alpha,\SuperQ_\beta\}=0,\qquad
    \{\SuperD_\alpha,\bar{\SuperD}_\beta\}
      =-2\gamma^\mu_{\alpha\beta}\partial_\mu .
  \]
  Hence \((\bar\alpha\SuperQ)\SuperD S=\SuperD(\bar\alpha\SuperQ)S\).
  A bare \(\theta\)-derivative lacks the compensating
  \(-\gamma^\mu\theta\partial_\mu\) term and does not commute with
  supersymmetry.

  \SupplementarySolution{5}{A Superspace Integration by Parts}
  The \(\theta\)-derivative part of \(\SuperD_\alpha S\) has Grassmann
  degree at most three and cannot contain a nonderivative \(D\)-component.
  Its remaining part is
  \(-(\gamma^\mu\theta)_\alpha\partial_\mu S\), whose \(D\)-component is a
  spacetime derivative.  Therefore
  \(\int d^4x[\SuperD_\alpha S]_D=0\).  Applying this to \(S_1S_2\) and
  using the graded Leibniz rule yields
  \[
    \int[S_1\SuperD_\alpha S_2]_D
      =-\int[(\SuperD_\alpha S_1)S_2]_D
      =-\int[S_2\SuperD_\alpha S_1]_D .
  \]

  \SupplementarySolution{6}{Solving the Chiral Constraint}
  Substitution in Eq.~\eqref{eq:26.2.26} gives
  \(\SuperD_Rx_+^\mu=0\) and \(\SuperD_R\theta_L=0\).
  A scalar solution is therefore an arbitrary polynomial in the two
  components of \(\theta_L\):
  \[
    \Phi=\phi(x_+)-\sqrt2\,
      (\theta_L^T\epsilon\psi_L(x_+))
      +\mathcal F(x_+)(\theta_L^T\epsilon\theta_L).
  \]
  This is Eq.~\eqref{eq:26.3.21}.  The complex \(\phi,\mathcal F\) give
  four real bosonic components, and the Majorana field represented by
  \(\psi_L\) gives four real fermionic components.

  \SupplementarySolution{7}{Multiplying Chiral Superfields}
  Direct multiplication in \(x_+,\theta_L\) coordinates gives
  \[
  \begin{aligned}
    \phi_{12}&=\phi_1\phi_2,\\
    \psi_{12,L}&=\phi_1\psi_{2L}+\phi_2\psi_{1L},\\
    \mathcal F_{12}
      &=\phi_1\mathcal F_2+\phi_2\mathcal F_1
        -(\psi_{1L}^T\epsilon\psi_{2L}).
  \end{aligned}
  \]
  Although the spinors anticommute, \(\epsilon\) is antisymmetric too, so
  \(\psi_{1L}^T\epsilon\psi_{2L}
   =\psi_{2L}^T\epsilon\psi_{1L}\).  The last line is consequently
  symmetric under interchange of the multiplets.

  \SupplementarySolution{8}{Why an \(\mathcal F\)-Term Is Invariant}
  Eq.~\eqref{eq:26.3.16} gives immediately
  \[
    \delta\mathcal F
      =\sqrt2\,\bar\alpha_L\sl{\partial}\psi_L
      =\partial_\mu(\sqrt2\,\bar\alpha_L\gamma^\mu\psi_L).
  \]
  Its spacetime integral is invariant for the usual boundary conditions.
  Likewise Eq.~\eqref{eq:26.2.17} gives
  \(\delta D=\partial_\mu(i\bar\alpha\gamma_5\gamma^\mu\lambda)\).
  Thus both top components change only by total derivatives.

  \SupplementarySolution{9}{Kahler Transformations}
  Both \(h(\Phi)\) and \(h(\Phi)^*\) are chiral superfields.  Chiral
  superfields have \(D\)-components that are spacetime derivatives, so
  \[
    \int d^4x\,[h(\Phi)+h(\Phi)^*]_D=0 .
  \]
  Also
  \(\partial_n\partial_{\bar m}h=
    \partial_n\partial_{\bar m}h^*=0\).  Hence
  \(K_{n\bar m}\), and therefore the component kinetic terms, are
  unchanged.

  \SupplementarySolution{10}{Renormalizability by Dimension Counting}
  Since an \(\mathcal F\)-component raises dimension by one and a
  \(D\)-component by two, a dimension-four Lagrangian requires
  \(d(f)\le3\) and \(d(K)\le2\).  With \(d(\Phi)=1\), the derivative-free
  possibilities are a cubic holomorphic polynomial for \(f\) and a
  quadratic function for \(K\).  Purely holomorphic or antiholomorphic
  terms in \(K\) have vanishing integrated \(D\)-terms.  Positivity and
  reality then leave
  \[
    K=g_{n\bar m}\Phi_n\Phi_m^*,\qquad
    g_{n\bar m}=g_{m\bar n}^*,
  \]
  while derivative terms are either too large in dimension or removable
  by the integrations described in Section 26.3.

  \SupplementarySolution{11}{Eliminating the Chiral Auxiliary Field}
  The auxiliary part is
  \[
    \mathcal L_{\rm aux}
      =\sum_n\left(\mathcal F_n^*\mathcal F_n
       \mathcal F_nf_n+\mathcal F_n^*f_n^*\right)
      =\sum_n|\mathcal F_n+f_n^*|^2-\sum_n|f_n|^2 .
  \]
  Therefore \(\mathcal F_n=-f_n^*\), and the on-shell Lagrangian contains
  \(-V\), with
  \[
    V(\phi,\phi^*)=\sum_n|f_n(\phi)|^2.
  \]

  \SupplementarySolution{12}{When Does an \(F\)-Term Vacuum Exist?}
  The potential is a sum of nonnegative terms.  If all \(f_n\) vanish at a
  point, then \(V=0\), the auxiliary expectation values vanish, and the
  supersymmetry variations of the fermions have no constant term.
  Conversely a supersymmetric vacuum is annihilated by the supercharges and
  has zero energy.  Positivity of
  \(V=\sum|f_n|^2\) then forces every \(f_n\) separately to vanish.  Thus a
  common zero is necessary and sufficient at tree level.

  \SupplementarySolution{13}{Symmetry of the Fermion Mass Matrix}
  Around \(\phi=v+\varphi\),
  \[
    f=f(v)+f_n(v)\varphi_n+\frac12f_{nm}(v)\varphi_n\varphi_m+\cdots .
  \]
  Eq.~\eqref{eq:26.4.4} supplies
  \(-\tfrac12f_{nm}\bar\psi_n\psi_{mL}\).  Ordinary mixed partial
  derivatives commute, so \(f_{nm}=f_{mn}\), as also required by the
  symmetry of the Majorana chiral bilinear.  Complex conjugation gives
  \(-\tfrac12f_{nm}^*\bar\psi_n\psi_{mR}\).

  \SupplementarySolution{14}{A Moduli-Space Tangent Vector}
  Differentiating the vacuum equations gives
  \[
    0=\frac{d f_n}{ds}
      =f_{nm}(\phi(s))\frac{d\phi_m}{ds}.
  \]
  Hence the tangent \(v_m=d\phi_m/ds\) is a zero eigenvector of the
  fermion mass matrix.  Supersymmetry pairs this massless fermion with the
  two real scalar fluctuations tangent to the complex vacuum family, giving
  one massless chiral multiplet.

  \SupplementarySolution{15}{The Goldstino Null Eigenvector}
  Stationarity with respect to \(\phi_m\) gives
  \[
    0=\frac{\partial V}{\partial\phi_m}
      =f_{nm}f_n^*.
  \]
  Thus \(f_n^*\) is a right null vector of the symmetric fermion mass
  matrix.  If \(F^2=\sum_n|f_n|^2>0\), the normalized field
  \[
    \psi_G=\frac{f_n^*\psi_{nL}}{F}
  \]
  is massless.  Its supersymmetry variation has the nonzero constant
  \(\delta\psi_G\propto F\alpha_L\), identifying it as the goldstino.

  \SupplementarySolution{16}{A Minimal O'Raifeartaigh Model}
  The derivatives are
  \[
    f_X=\lambda\Phi^2-\mu^2,\qquad
    f_Y=m\Phi,\qquad
    f_\Phi=2\lambda X\Phi+mY.
  \]
  The first two cannot both vanish for \(m\mu\ne0\).  At
  \(\Phi=Y=0\), \(X\) is arbitrary and \(V=|\mu|^4\).
  Expanding in \(\Phi\) shows local stability, for example at \(X=0\),
  when \(|m|^2>2|\lambda\mu^2|\); the precise inequality follows from the
  two real \(\Phi\) eigenvalues.  The field \(X\) is the tree-level
  pseudomodulus and \(\psi_X\) is the goldstino.

  \SupplementarySolution{17}{The Current as a Moment Map}
  Infinitesimal invariance of \(K\) gives
  \[
    0=i\epsilon\left(
      K_n\mathcal T_{nm}\Phi_m
      -K_{\bar n}\mathcal T_{mn}\Phi_m^*\right).
  \]
  Therefore the two expressions for \(\mathcal J\) are equal and are
  complex conjugates, so \(\mathcal J\) is real.  Adding a real constant
  preserves this property and changes none of its superderivatives or its
  vector current, which is the usual moment-map ambiguity.

  \SupplementarySolution{18}{The Canonical \(U(1)\) Current Multiplet}
  The product formulas give
  \[
  \begin{aligned}
    C^{\mathcal J}&=|\phi|^2,\\
    \omega^{\mathcal J}&=-i\sqrt2\gamma_5
       (\phi^*\psi_L+\phi\psi_R),\\
    V_\mu^{\mathcal J}&=-i(\phi^*\partial_\mu\phi
       -\phi\partial_\mu\phi^*)
       +i\bar\psi_R\gamma_\mu\psi_L.
  \end{aligned}
  \]
  Varying the canonical scalar and fermion kinetic terms under
  \(\delta\phi=i\epsilon\phi,\ \delta\psi_L=i\epsilon\psi_L\) yields this
  same vector, up to the overall sign fixed by whether the current is
  defined from \(\delta\mathcal L=-(\partial_\mu\epsilon)J^\mu\) or with
  the opposite convention.

  \SupplementarySolution{19}{Components of a Linear Superfield}
  The independent coefficients of
  \(\SuperD_R^2S=\SuperD_L^2S=0\) give
  \[
    M=N=0,\qquad \partial_\mu V^\mu=0,\qquad
    \lambda=-\sl{\partial}\omega,\qquad D=-\Box C.
  \]
  Thus \(C,\omega\), and a divergence-free \(V_\mu\) determine the
  multiplet.  Off shell this is \(1+3=4\) bosonic components and four
  Majorana fermionic components; \(\lambda,D\) are derivative descendants,
  not independent auxiliaries.

  \SupplementarySolution{20}{The Superfield Euler--Lagrange Equation}
  With \(\delta\Phi_n=\SuperD_R^2\delta S_n\), superspace integration by
  parts moves both derivatives onto \(K_n\):
  \[
    \delta I_K\propto
      \int d^4x\,d^4\theta\,
      \delta S_n\SuperD_R^2K_n .
  \]
  The chiral projection identity
  \(\int[(\SuperD_R^2h)]_{\mathcal F}
    =2\int[h]_D\), together with the normalizations
  \eqref{eq:26.6.5}--\eqref{eq:26.6.6}, turns the superpotential variation
  into \(4f_n\) in the same equation.  Arbitrary \(\delta S_n\) therefore
  gives
  \[
    \SuperD_R^2K_n=-4f_n,
  \]
  as in Eq.~\eqref{eq:26.6.15}.

  \SupplementarySolution{21}{The Kahler Metric in Components}
  The mixed quadratic term in the \(\theta\)-expansion of \(K\) yields
  \[
    \mathcal L_{\rm scalar}
      =-K_{n\bar m}(\phi,\phi^*)
        \partial_\mu\phi_n\partial^\mu\phi_m^* .
  \]
  Around any constant background, diagonalize the Hermitian matrix
  \(K_{n\bar m}\).  A negative eigenvalue gives a scalar mode with a
  wrong-sign kinetic term and hence negative norm or unbounded energy.
  Absence of such ghosts requires, and locally is guaranteed by, positive
  definiteness of the Kahler metric.

  \SupplementarySolution{22}{Auxiliaries with Noncanonical \(K\)}
  Varying Eq.~\eqref{eq:26.8.6} with respect to
  \(\mathcal F_{\bar n}^*\) gives
  \[
    K_{m\bar n}\mathcal F_m
      -\frac12K_{mp\bar n}(\bar\psi_m\psi_{pL})
      +f_{\bar n}^*=0.
  \]
  Its conjugate determines \(\mathcal F_m^*\).  With fermions set to zero,
  \(\mathcal F_m=-(K^{-1})^{m\bar n}f_{\bar n}^*\).  Completing the
  square then yields
  \[
    V=f_n(K^{-1})^{n\bar m}f_{\bar m}^*.
  \]

  \SupplementarySolution{23}{Kahler Curvature and Four Fermions}
  Solving the auxiliary equation produces a term quadratic in the Kahler
  connection \(K_{np\bar q}\).  Combining it with the explicit
  \(K_{np\bar m\bar q}\) term gives
  \[
    R_{n\bar m p\bar q}
      =K_{n\bar m p\bar q}
       -K_{np\bar r}(K^{-1})^{s\bar r}K_{s\bar m\bar q}.
  \]
  Hence the four-fermion interaction is proportional to
  \(R_{n\bar m p\bar q}
  (\bar\psi_n\psi_{pL})(\bar\psi_q\psi_{mR})\).
  Reordering the Grassmann fields can move an overall sign or a factor
  between this formula and the definition of \(R\), but the covariant
  curvature tensor is invariant.

  \SupplementarySolution{24}{Conservation of the Supersymmetry Current}
  Differentiating Eq.~\eqref{eq:26.7.10} produces terms containing
  \(\Box\phi_n\), \(\sl{\partial}\psi_{nL,R}\), and
  \(\partial_\mu f_n=f_{nm}\partial_\mu\phi_m\).
  Use Eq.~\eqref{eq:26.7.17} for the fermions and the scalar equation
  \[
    \Box\phi_n=f_{nm}^*f_m
      +\frac12f_{nmp}^*(\bar\psi_m\psi_{pR})
  \]
  with its conjugate.  The \(f_{nm}\partial\phi\) terms cancel the
  differentiated superpotential terms, while the remaining Yukawa terms
  cancel pairwise by \(f_{nm}=f_{mn}\).  Thus
  \(\partial_\mu S^\mu=0\).

  \SupplementarySolution{25}{An Identically Conserved Improvement}
  The derivatives commute and are symmetric in \(\mu,\nu\), whereas
  \([\gamma^\mu,\gamma^\nu]\) is antisymmetric:
  \[
    \partial_\mu\partial_\nu
      ([\gamma^\mu,\gamma^\nu]\chi)=0.
  \]
  The time component of the improvement is a spatial divergence because
  the \(\nu=0\) term vanishes.  Its integral is therefore a boundary term,
  so it changes neither the charge nor the commutators of that charge with
  fields.

  \SupplementarySolution{26}{Gamma Trace and Scale Invariance}
  For homogeneous cubic \(f\), Euler's theorem gives
  \(\phi_mf_m=3f\).  Differentiating with respect to \(\phi_n\) yields
  \[
    \phi_m f_{nm}=2f_n,
  \]
  so every parenthesis in Eq.~\eqref{eq:26.7.19} vanishes.  A quadratic
  term instead gives \(\phi_mf_{nm}=f_n\), and a linear term gives zero;
  both leave a nonzero difference from \(2f_n\).  These dimensionful mass
  and tadpole couplings are precisely the classical obstructions to a
  gamma-traceless improved current.

  \SupplementarySolution{27}{Inverting the Supercurrent Relation}
  Gamma tracing gives
  \[
    \gamma_\mu S^\mu
      =-2\gamma_\mu\omega^\mu
       +2\gamma_\mu\gamma^\mu\gamma^\nu\omega_\nu
      =6\gamma^\nu\omega_\nu.
  \]
  Thus
  \[
    \omega_\mu=-\frac12S_\mu
      +\frac16\gamma_\mu\gamma^\nu S_\nu.
  \]
  Substitution back into the original relation verifies the inverse.

  \SupplementarySolution{28}{Parity Exchanges Chiralities}
  The parity substitution gives
  \(x_+\mapsto\Lambda_Px_-\) and
  \(P_L(-i\beta\theta)=-i\beta P_R\theta\).  Hence a function only of
  \(x_+,\theta_L\) becomes one only of \(x_-,\theta_R\), so left chirality
  becomes right chirality.  Up to the intrinsic phase,
  \[
    \phi(x)\mapsto\phi(\Lambda_Px),\quad
    \psi_L(x)\mapsto-i\beta\psi_L(\Lambda_Px)\ \hbox{in the right slot},
    \quad
    \mathcal F(x)\mapsto-\mathcal F(\Lambda_Px),
  \]
  consistently with Eqs.~\eqref{eq:26.3.35}--\eqref{eq:26.3.36}.

  \SupplementarySolution{29}{Time Reversal Preserves Chirality}
  The matrix \(B_T\) intertwines complex-conjugate gamma matrices with the
  time-reflected ones.  When \(T\) is moved through
  \(\SuperD_R\Phi=0\), its antiunitarity also sends \(i\to-i\).  This
  extra sign compensates the interchange that occurred for unitary parity,
  giving
  \[
    \SuperD_R\Phi^T=0.
  \]
  Equivalently, time reversal flips both momentum and spin, so helicity is
  unchanged; parity flips momentum but not angular momentum, so helicity
  and chirality are exchanged.

  \SupplementarySolution{30}{A Dimension-Five Operator Basis}
  An \(\mathcal F\)-integrand must have dimension four and a \(D\)-integrand
  dimension three.  The nonderivative representatives are
  \[
    [\Phi^4]_{\mathcal F}+\text{c.c.},\qquad
    [\Phi^2\Phi^*]_D+\text{c.c.}
  \]
  Purely chiral cubic \(D\)-terms vanish.  The only independent
  higher-derivative scalar class may be represented by
  \[
    [\Phi\Box\Phi]_{\mathcal F}+\text{c.c.}
  \]
  or, after chiral projection and superspace integration by parts, by the
  corresponding \(D\)-term with two spinor superderivatives.  Lorentz
  invariance excludes one ordinary derivative, and all other placements of
  two spinor derivatives reduce to this representative.  These three
  classes exhaust the basis.
\end{enumerate}
